
\begin{abstract}
Computerized Adaptive Testing (CAT), as a key technology for realizing personalized education, aims to accurately assess an examinee's proficiency by retrieving exercises from a candidate pool that dynamically match their current ability estimates. However, existing CAT research paradigms have long been constrained by the dual limitations of static offline data and isolated component optimization. Severely restricted by the issue of partial labels in offline interaction logs, researchers are forced to degrade the inherently dynamic assessment process into a static sequence prediction problem. Meanwhile, current research predominantly focuses on isolated perspectives—such as either the selection strategy or the diagnostic model—neglecting the holistic investigation of the overall CAT interaction process. To break through these bottlenecks, this paper proposes AgentCAT, a Large Language Model-based multi-agent simulation system, aiming to construct a high-fidelity benchmarking simulation environment for dynamic testing. This framework comprises three core modules: (1) The examinee agent, equipped with memory retrieval and Chain-of-Thought reasoning capabilities, acts as a cognitive simulator to replicate authentic response behaviors based on cognitive profiles; (2) The selection agent acts as an instructional planner, employing a coarse-to-fine bucketing strategy and a knowledge graph exploration mechanism to strike a balance between local difficulty matching and global knowledge coverage; (3) The supervisor introduces a dual-auditing and robust update mechanism to ensure the convergence and statistical validity of ability estimation. To comprehensively validate the effectiveness of the framework, we conducted systematic evaluations on two real-world datasets across three dimensions: macro-level ability convergence, micro-level instructional interaction logic, and the capability to overcome data sparsity. Experimental results demonstrate that AgentCAT not only achieves highly effective ability estimation, but its selection strategy also successfully strikes a balance between difficulty adaptation and instructional coherence that perfectly aligns with human pedagogical intuition.
\end{abstract}